% ---
% Síntesis global por objetivo específico (OE-1 a OE-6) y evaluación por dimensiones.
% Sección 6.8 del reporte técnico.
% Entradas que lo respaldan: H-001, H-005, H-020, H-023, H-024, H-026, H-035, H-038, H-039, H-040, H-041, H-042, H-043, H-049, H-051, H-052, H-032, D-012, D-014, D-016, D-031, D-036.
% ---

\subsection{Síntesis por objetivo específico}
\label{sec:sintesis_objetivos}

% fuente: H-039, H-040, H-042, H-044; 6_alineamiento_kernel.csv, 6_dimension_efectiva.csv, 6_diferencia_geometrica_resumen.csv,
% 6_metricas_embedding.csv, 7_cv_resumen.csv, 7_test.csv, X4_test.csv, 5a_resumen.csv, 5a_tiempos_clasicos.csv.
% Tabla añadida el 2026-10-20; tiempos de C2 y C3 medidos con Code/tools/tiempos_transformacion.py (5a no los cronometraba).
La Tabla~\ref{tab:resumen_resultados} reúne en un solo lugar las métricas principales de las cinco condiciones.
En la familia de kernels, el exceso de alineamiento de C4 y C5 (0.003 a 0.007) queda muy por debajo del kernel RBF
(0.049 en masas y 0.151 en calcificaciones) y su dimensión efectiva, cercana a 170--190 de 200, refleja la concentración.
El mismo MLP clasifica peor sobre C4 y C5 que sobre C1--C3, mientras que una máquina de vectores de soporte sobre
el kernel de C4 queda cerca de la que usa el kernel RBF. En costo, PCA y Kernel PCA transforman todas las lesiones
en menos de 0.2\,s, frente a 13 a 16\,s de los dos mapas cuánticos simulados con vectores de estado.

% Generado por Code/tools/tablas_latex.py desde 6_alineamiento_kernel.csv, 6_dimension_efectiva.csv, 6_diferencia_geometrica_resumen.csv, 6_metricas_embedding.csv, 7_cv_resumen.csv, 7_test.csv, X4_test.csv, 5a_resumen.csv y 5a_tiempos_clasicos.csv. No editar a mano: volver a correr el script.
\begin{table}[H]
\centering
\small
\caption{Resumen de resultados por condición: separabilidad (M6), clasificación (M7 y X4) y tiempo de transformación (M5).}
\label{tab:resumen_resultados}
\begin{adjustbox}{max width=\linewidth}
\begin{tabular}{l r r r r r r c r}
\toprule
 & \multicolumn{3}{c}{Kernel} & \multicolumn{2}{c}{Representación} & \multicolumn{2}{c}{AUC-ROC} & \\ \cmidrule(lr){2-4}\cmidrule(lr){5-6}\cmidrule(lr){7-8}
Condición & Exceso KTA $\uparrow$ & $d_{\mathrm{eff}}$ & $g$ & DB $\downarrow$ & Fisher $\uparrow$ & VC & Prueba [IC 95\,\%] & Tiempo (s) \\
\midrule
\multicolumn{9}{l}{\textit{Masas}} \\
\quad C1 & — & — & — & \textbf{4.301} & 0.669 & \textbf{0.652} & \textbf{0.651} [0.578, 0.724] & — \\
\quad C2 & — & — & — & \textbf{4.301} & 0.669 & 0.641 & 0.641 [0.572, 0.710] & 0.0002 \\
\quad C3 & \textbf{0.049} & 2.8 & — & 5.564 & \textbf{0.781} & 0.648 & 0.646 [0.572, 0.721] & 0.14 \\
\quad C4 & 0.003 & 184.0 & 1.433 & 5.898 & 0.449 & 0.565 & 0.530 [0.473, 0.587] & 13.3 \\
\quad C5 & 0.005 & 171.3 & 1.511 & 13.346 & 0.668 & 0.563 & 0.581 [0.521, 0.644] & 14.3 \\
\quad SVM sobre el kernel RBF & \multicolumn{6}{c}{} & 0.649 [0.574, 0.726] & — \\
\quad QSVM sobre el kernel de C4 & \multicolumn{6}{c}{} & 0.628 [0.555, 0.703] & — \\
\quad QSVM sobre el kernel de C5 & \multicolumn{6}{c}{} & 0.599 [0.523, 0.675] & — \\
\multicolumn{9}{l}{\textit{Calcificaciones}} \\
\quad C1 & — & — & — & \textbf{2.387} & 1.284 & 0.742 & 0.769 [0.694, 0.836] & — \\
\quad C2 & — & — & — & \textbf{2.387} & 1.284 & \textbf{0.758} & \textbf{0.788} [0.719, 0.851] & 0.0002 \\
\quad C3 & \textbf{0.151} & 2.6 & — & 2.995 & \textbf{1.390} & 0.739 & 0.758 [0.685, 0.826] & 0.20 \\
\quad C4 & 0.006 & 187.7 & 1.400 & 6.890 & 0.399 & 0.640 & 0.633 [0.558, 0.700] & 13.9 \\
\quad C5 & 0.007 & 167.6 & 1.701 & 6.247 & 0.753 & 0.667 & 0.652 [0.571, 0.725] & 16.1 \\
\quad SVM sobre el kernel RBF & \multicolumn{6}{c}{} & 0.769 [0.696, 0.834] & — \\
\quad QSVM sobre el kernel de C4 & \multicolumn{6}{c}{} & 0.759 [0.687, 0.827] & — \\
\quad QSVM sobre el kernel de C5 & \multicolumn{6}{c}{} & 0.701 [0.620, 0.775] & — \\
\bottomrule
\end{tabular}
\end{adjustbox}
\par\smallskip
\parbox{0.95\linewidth}{\footnotesize Separabilidad sobre 200 lesiones de entrenamiento por subconjunto y $c = 1$. Las columnas de kernel de C3 corresponden al kernel RBF, con el mismo $\gamma$ que C3; C1 y C2 no definen un kernel propio. $g$: mínimo de la \textit{geometric difference} regularizada ($\sqrt{N} \approx 14$). AUC: el mismo MLP en C1--C5; las tres últimas filas de cada bloque son clasificadores de vectores de soporte sobre el kernel (experimento X4). Tiempo: transformación de todas las lesiones del subconjunto en la máquina del proyecto; en C2 y C3, mediana de cinco repeticiones; C1 es la identidad. Con el etiquetado diseñado según Huang et al., control positivo del protocolo (sección~\ref{sec:qsvm_ventaja}), el kernel cuántico proyectado alcanza en validación cruzada 0.998 en masas y 0.994 en calcificaciones, frente a 0.800 y 0.824 del mejor clasificador clásico. En negritas, el mejor valor de cada columna con dirección en cada subconjunto; en el AUC de prueba se comparan también las filas de X4.}
\end{table}


% Añadido el 2026-10-23: vista de conjunto de todas las comparaciones del trabajo, una por experimento.
La Tabla~\ref{tab:mapa_comparaciones} ordena por experimento todas las comparaciones del trabajo: qué
pregunta respondió cada una, qué se comparó, el resultado principal y su lectura. En la tarea diagnóstica,
lo clásico va por delante o, en el mejor caso para lo cuántico, empata; lo cuántico solo gana con un
etiquetado diseñado a partir de su propia geometría.

% fuente: H-039, H-042, H-043, H-035, H-044, H-047, H-053, H-054, H-056, H-050, H-048, H-057; 6_alineamiento_kernel.csv, 7_test.csv, X1_auc_por_condicion.csv, X4_test.csv, X6_test.csv, X7_resumen.csv, X8_test.csv, X10_test.csv, X11_transicion.csv, X6_etiquetas_construidas.csv
\begin{table}[htbp]
\centering
\small
\caption{Mapa de las comparaciones del trabajo. AUC-ROC en el conjunto de prueba salvo que se indique validación cruzada (VC); masas / calcificaciones. En negritas, el mejor valor de cada comparación.}
\label{tab:mapa_comparaciones}
\begin{adjustbox}{max width=\linewidth}
\begin{tabular}{p{0.1\linewidth} p{0.2\linewidth} p{0.2\linewidth} p{0.27\linewidth} p{0.17\linewidth}}
\toprule
\textbf{Exp.} & \textbf{Pregunta} & \textbf{Comparación} & \textbf{Resultado} & \textbf{Lectura} \\
\midrule
M6 (\S\ref{sec:separabilidad_representaciones}) & ¿Separan mejor las clases las representaciones cuánticas? & Métricas de kernel y de representación en C1--C5 & Exceso de KTA de 0.003 a 0.007 en C4 y C5, frente a \textbf{0.049} / \textbf{0.151} del RBF & Clásicas por delante en todas las métricas \\
M7 (\S\ref{sec:rendimiento_clasificacion}) & ¿Clasifica mejor el mismo perceptrón sobre C4 o C5? & MLP $12 \to 32 \to 16 \to 1$ en C1--C5 & C1: \textbf{0.651} / \textbf{0.769}; C4: 0.530 / 0.633; C5: 0.581 / 0.652 & Clásicas, de 0.07 a 0.14 por delante \\
M7, sensibilidad & ¿Depende de la red, la semilla, \texttt{reps} o $k$? & Semillas, redes 16-8 y 64-32, sin estandarizar, 200 épocas, \texttt{reps} = 2, $k = 8$ & La semilla mueve el AUC como mucho 0.039; la brecha se mantiene & Sin cambio \\
X1 & ¿Cuánto cuesta usar 12 de las 67 características? & Regresión logística con 67 y con 12 & \textbf{0.693} / \textbf{0.800} frente a 0.650 / 0.758 (VC) & Pérdida de 2 a 4 puntos, igual para todas las condiciones \\
X4 & ¿Se pierde la información al codificar o al leer? & QSVM de C4, SVM RBF y MLP sobre $\langle Z_i\rangle$ & 0.628 / 0.759; \textbf{0.649} / \textbf{0.769}; 0.530 / 0.633 & Empate entre QSVM y SVM; la pérdida está en la lectura \\
X6 & ¿Gana algún kernel cuántico con la misma búsqueda en ambos lados? & 28 candidatos cuánticos frente a 12 clásicos, búsqueda anidada & \textbf{0.656} frente a 0.652 / 0.755 frente a \textbf{0.760} & Empate \\
X7 & ¿Sesgó el sorteo de $\boldsymbol\theta$ el resultado? & Diez sorteos de $\boldsymbol\theta$ & Ninguno alcanza a C1; déficit de 0.03 a 0.18 (VC) & Clásicas \\
X8 & ¿Ayuda entrenar $\boldsymbol\theta$? & $\boldsymbol\theta$ entrenado frente a fijo & Gana de $+0.003$ a $+0.041$; sigue entre $-0.066$ y $-0.118$ frente a C1 & Clásicas \\
X10 & ¿Cambia con árboles o con tres \textit{folds}? & Árboles, \textit{random forest} y XGBoost en C1--C5; MLP con tres \textit{folds} & \textit{Random forest}, C1: \textbf{0.636} / \textbf{0.768}; C4: 0.591 / 0.651; C5: 0.614 / 0.624 & Clásicas; en masas, empate con los ensambles \\
M8 (\S\ref{sec:comparacion_global}) & ¿Cuánto separan los descriptores del radiólogo? & Regresión logística sobre forma y margen, o tipo y distribución & \textbf{0.872} / \textbf{0.825} (VC) & Referencia clínica, no competidor \\
X6, X11 & ¿Existe una tarea con ventaja cuántica sobre estas lesiones? & QSVM proyectado frente al mejor clásico con el etiquetado diseñado & \textbf{0.998} frente a 0.857 / \textbf{0.994} frente a 0.876 (VC); $\alpha^{*}$ de 0.4 a 0.7 & Cuántico, solo con el etiquetado diseñado \\
\bottomrule
\end{tabular}
\end{adjustbox}
\end{table}

A fin de articular los hallazgos empíricos obtenidos a lo largo del protocolo experimental con los compromisos
planteados al inicio de la investigación, la Tabla~\ref{tab:sintesis_objetivos_especificos} presenta una síntesis
estructurada por cada uno de los seis objetivos específicos (OE-1 a OE-6) formulados en el planteamiento del trabajo.
Para cada objetivo se resume el alcance establecido, la evidencia cuantitativa central generada, la sección del
reporte donde se desarrolla el análisis detallado y la respuesta concreta alcanzada en función de los datos.

% fuente: H-001, H-005, H-020, H-023, H-024, H-026, H-035, H-038, H-039, H-040, H-041, H-042, H-043, H-049, H-051, H-052, H-032, D-012, D-014, D-016, D-031, D-036; 2_preprocesamiento_qa.csv, X1_varianza_retenida.csv, B1_benchmark_arquitectura.csv, 6_alineamiento_kernel.csv, 6_dimension_efectiva.csv, 6_diferencia_geometrica_resumen.csv, 7_cv_por_fold.csv, 7_test.csv, 7_test_diferencias.csv, 8_spearman.csv, 8_oe6_qubits.csv, 8_oe6_shots_kernel.csv, 8_oe6_shots_embedding.csv, 8_oe6_resumen.csv
\begin{table}[htbp]
\centering
\small
\caption{Síntesis de resultados y evidencia empírica por objetivo específico.}
\label{tab:sintesis_objetivos_especificos}
\begin{adjustbox}{max width=\linewidth}
\begin{tabular}{p{0.07\linewidth} p{0.24\linewidth} p{0.31\linewidth} p{0.14\linewidth} p{0.24\linewidth}}
\toprule
\textbf{Objetivo} & \textbf{Objetivo resumido} & \textbf{Evidencia principal} & \textbf{Sección} & \textbf{Resultado} \\
\midrule
OE-1 & Construir pipeline sobre CBIS-DDSM (DICOM, ROI, 67 variables radiómicas, selección a $k = 12$ y escalado angular $[0, \pi]$ en subproblemas independientes). & 3\,568 registros con cobertura DICOM 100\,\%; 67 variables extraídas; 12 seleccionadas retienen 86.8\,\% (masas) y 90.2\,\% (calcificaciones) de varianza; 0 solapes entre train y test. & Sección~\ref{sec:datos_caracteristicas_seleccion} & Pipeline validado sin fuga de pacientes; selección por prueba F con filtro de redundancia ($|r| \le 0.95$) y escalado por cuantiles con IQR de $\pi/2$. \\
\midrule
OE-2 & Implementar cinco condiciones experimentales de dimensión $k = 12$ (C1 crudo, C2 PCA lineal, C3 Kernel PCA RBF, C4 ZZFeatureMap, C5 PauliFeatureMap). & 5 condiciones simétricas sobre el mismo vector; $\boldsymbol\theta$ fijo (semilla 42) en Arquitectura~B; C5 diferenciado con compuertas $Y$ y $ZZ$; extracción analítica en 0.20\,h frente a 381.7\,h estimadas en el benchmark de recursos para entrenar los dos mapas con Arquitectura~A. & Secciones~\ref{sec:modulo5}, \ref{sec:lecciones_implementacion} & Condiciones implementadas; control nulo C2 preservado; Arquitectura~B desacopla codificación de clasificación favoreciendo la reproducibilidad. \\
\midrule
OE-3 & Analizar y comparar separabilidad de clases en las cinco condiciones mediante KTA, diferencia geométrica $g$, Davies--Bouldin y razón de Fisher. & A $c = 1$, exceso de KTA de 0.003 a 0.007 en C4/C5 frente a 0.049 y 0.151 en RBF; $d_{\mathrm{eff}}$ cuántica de 167.6 a 187.7; $g$ mínima de 1.40 a 1.70 ($\sqrt{N} \approx 14.1$); condiciones clásicas superan a cuánticas en Fisher. & Sección~\ref{sec:separabilidad_representaciones} & Las transformaciones cuánticas a $c = 1$ no muestran mayor separabilidad que las clásicas en ninguna familia métrica; matrices de kernel concentradas próximas a la identidad. \\
\midrule
OE-4 & Entrenar un mismo clasificador MLP sobre las cinco representaciones bajo idéntica arquitectura, hiperparámetros y validación cruzada 5-fold por paciente. & Red $12 \to 32 \to 16 \to 1$, optimizador Adam, parada temprana; variación máxima por semilla de 0.039 en CV; 200 épocas sin parada temprana no mejora C4/C5 (AUC 0.533 y 0.557 en masas; 0.596 y 0.631 en calcificaciones). & Sección~\ref{sec:rendimiento_clasificacion} & Protocolo simétrico validado; convergencia uniforme; el déficit cuántico no es atribuible a subentrenamiento ni a sensibilidad de hiperparámetros. \\
\midrule
OE-5 & Evaluar si la separabilidad se traslada a clasificación binaria (AUC-ROC, F1, Accuracy, Precision, Recall) y correlación de Spearman entre rankings. & En test, AUC masas: C1 0.651, C3 0.646, C4 0.530, C5 0.581; calcificaciones: C1 0.769, C3 0.758, C4 0.633, C5 0.652. Diferencias frente a C1 entre $-0.07$ y $-0.14$. Coeficiente $\rho_s$ descriptivo en el rango de 0.67 a 0.97. & Secciones~\ref{sec:rendimiento_clasificacion}, \ref{sec:comparacion_global} & El MLP clasifica peor sobre C4/C5; separabilidad y clasificación coinciden en su ordenamiento (C1--C3 delante de C4--C5), pero sin ganancia previa en separabilidad que trasladar. \\
\midrule
OE-6 & Evaluar restricciones prácticas y del simulador cuántico (recursos, escalabilidad en qubits, gradientes y shots) para determinar viabilidad en salud. & Simulación exacta viable hasta $k \approx 25$ en 24\,h; gradiente por \textit{parameter-shift} requiere 5.16\,s/muestra ($\sim 205$\,h a $k = 12$); a 1\,024 disparos, entre 15.7\,\% y 74.5\,\% de pares del kernel leen cero; MLP tolera disparos perdiendo $< 0.02$ de AUC. & Secciones~\ref{sec:restricciones_simulador}, \ref{sec:lecciones_implementacion} & Techo práctico impuesto por gradientes físicos y concentración del kernel; fidelidad experimentalmente inmedible con disparos realistas en NISQ; atajos de simulación no trasladables a hardware. \\
\bottomrule
\end{tabular}
\end{adjustbox}
\end{table}

A partir de esta correspondencia empírica, la respuesta integral a la pregunta de investigación se estructura
en torno a las tres dimensiones de análisis que fundamentaron el planteamiento del problema:

% fuente: H-038, H-039, H-040, H-041; 6_alineamiento_kernel.csv, 6_dimension_efectiva.csv, 6_diferencia_geometrica_resumen.csv
\textbf{Dimensión de representación (OE-3).} Respecto a la capacidad de las transformaciones cuánticas
variacionales para generar espacios de características con mayor separabilidad de clases, la evidencia empírica
indica de forma consistente que los mapas cuánticos evaluados a escala nominal ($c = 1$) no superan a las
representaciones clásicas de referencia en ninguna de las familias métricas analizadas. En el espacio de kernels,
el exceso de alineamiento centrado (\textit{kernel-target alignment}, KTA) de C4 y C5 se sitúa entre 0.003 y 0.007,
abarcando menos del 10\,\% del exceso registrado por el kernel gaussiano RBF clásico (0.049 en masas y 0.151 en
calcificaciones); las matrices cuánticas resultan cercanas a la identidad con una dimensión efectiva que satura
el espacio disponible ($d_{\mathrm{eff}} \approx 168$--$188$ sobre 200), y la diferencia geométrica regularizada
($g \approx 1.40$--$1.70$) permanece muy alejada de la cota teórica de ventaja cuántica potencial ($\sqrt{N} \approx 14.1$).
Asimismo, en las representaciones locales $\langle Z_i \rangle$, los índices de Fisher y Davies--Bouldin sitúan a
las condiciones clásicas C1 y C3 sistemáticamente por delante de las condiciones cuánticas. Aunque la contracción
angular ($c \to 0.02$) o la reducción a 8 qubits (X3) atenúan la concentración, dicha recuperación no introduce
estructura cuántica genuina: el alineamiento solo se aproxima al clásico en el régimen en que el producto cruzado
entre características deja de influir, emulando la geometría de un kernel clásico.

% fuente: H-042, H-043, H-049; 7_test.csv, 7_test_diferencias.csv, 8_spearman.csv
\textbf{Dimensión de clasificación (OE-5).} En lo concerniente a si la transformación cuántica se traduce en una
mejora en la clasificación binaria de lesiones mamográficas, los resultados sobre el clasificador perceptrón multicapa
muestran que el mismo modelo clasifica de manera desfavorable sobre las representaciones cuánticas C4 y C5 en comparación
con las condiciones clásicas. En el conjunto de prueba independiente, las condiciones cuánticas exhiben un déficit de
0.07 a 0.14 de AUC-ROC frente al vector de características crudas C1 (alcanzando 0.530 en C4 masas y 0.633 en C4 calcificaciones,
frente a 0.651 y 0.769 en C1), brecha estadísticamente significativa respaldada por intervalos de confianza por remuestreo
que excluyen el cero. Este menoscabo no responde a artefactos de optimización, pues los análisis de sensibilidad sugieren
que ni la variación de semillas, ni arquitecturas alternativas, ni el entrenamiento prolongado a 200 épocas sin parada temprana
logran revertir la diferencia. El análisis de correlación de rangos de Spearman ($\rho_s \in [0.67, 0.97]$) muestra una coincidencia
descriptiva en el ordenamiento relativo de ambas etapas (con el bloque clásico C1--C3 situado por delante de C4 y C5); no obstante, debido a que
las condiciones cuánticas no alcanzaron mayor separabilidad en OE-3, no existió una ganancia representacional que pudiera
trasladarse al desempeño diagnóstico.

% fuente: H-051, H-052, H-054, H-019, H-032; 8_oe6_qubits.csv, 8_oe6_shots_kernel.csv, 8_oe6_resumen.csv, X8_costo_gradiente.csv
\textbf{Dimensión de viabilidad práctica (OE-6).} En relación con la factibilidad computacional y operacional de trasladar
estos esquemas a un contexto clínico real, las mediciones revelan restricciones fundamentales impuestas por el formalismo
físico de la medición cuántica. Si bien el clasificador de red neuronal tolera el ruido estocástico inducido por un régimen
realista de 1\,024 disparos por muestra sobre las expectativas locales $\langle Z_i \rangle$ (reteniendo el AUC exacto a menos
de 0.02), la estimación de matrices de kernel de fidelidad en procesadores cuánticos físicos resulta experimentalmente inviable
a escala nominal ($c = 1$): debido a la concentración exponencial de la fidelidad cuántica, entre el 15.7\,\% y el 74.5\,\%
de los pares de lesiones registrarían lecturas idénticamente nulas a 1\,024 disparos, y acotar el error relativo al 10\,\% exigiría una mediana de decenas de miles de
mediciones por par en C4 y de millones en C5. Asimismo, el coste computacional del gradiente exacto por \textit{parameter-shift}
(5.16\,s por muestra a $k = 12$, equivalente a unas 205~h de cálculo para 50 épocas) restringe de forma severa el entrenamiento
variacional en ausencia de los atajos de vector de estado analítico, los cuales solo existen en simulación clásica y son
inaplicables en dispositivos cuánticos físicos.
